%
%% Logic + proposed edits diagram:
%%   ../1-feedback/v0516/02-05_trait-rating-correlation_logic.txt
%% To open side-by-side: right-click the path above ->
%%   Open file -> "Open file under cursor in editor beside active"
%%   (or Alt+P for same-tab open; requires "Open file under cursor" extension by seito)
%

\subsection{Trait-Rating Correlation}

\begin{figure}[htbp]
\centering
\includegraphics[width=\textwidth]{0-display/Figures/node8/trait_satisfaction_violin_plots.pdf}
\caption{Trait scores and the review rating score. Area-scaled violin plots show review rating score distributions for discrete trait levels (0, 0.25, 0.5, 0.75, 1.0); violin width is proportional to sample size. The review rating score rises monotonically with trait score across all 10 traits, most steeply for patient-oriented subjective judgments (SPS, IQC, STS).}
\label{fig:trait_satisfaction}
\end{figure}


% =========================================================
% Para [trait-rating.setup] Setup -- does each trait retain variance independent of the rating?
% =========================================================
%% {CC-content-v0531}: [trait-rating.setup] a Setup para also delivers a result (S3) and a caveat (S4) -- 4 jobs. | keep S1-S2 as the setup/question; move the S3 result into P2. ========>
%
%% ---- P1.S1 [NEW v0519: hand-off fact from \S2.4] ----
In the previous subsection, we established that the ten traits remain distinguishable from each other.
%
%% ---- P1.S2 [NEW v0519: \S2.5 question, Test 2 of paired discriminant story] ----
We now ask whether each trait retains variance independent of review rating scores.
%
%% ---- P1.S3 [was P1.S1 in v0518; question-tail dropped, new P1.S2 carries it] ----
Trait--satisfaction correlations reach $r = 0.81$ (Figure~\ref{fig:trait_satisfaction}).
%% {CC-content-v0531}: r=0.81 here vs r=0.814 in P2.S2 for the same quantity. | reconcile in the values pass -- % TODO[values]. ========>
%% {CC-content-v0531}: "trait--satisfaction" names the same measure the caption calls "review rating score". | pick one term (flag for consistency pass). ========>
%
%% ---- P1.S4 [was P1.S2 in v0518; renumbered after S1/S2 inserts] ----
Such correlations are expected under a patient-perceived framing since both signals come from the same review narrative; the operative test is whether they survive controlling for review rating scores.
%% {CC-content-v0531}: two assertions joined by ";" (why expected; the operative test). | split into two sentences. ========>
%


% =========================================================
% Para [trait-rating.partial] Result -- correlations are strong but partial correlations stay non-zero
% =========================================================
%
%% ---- P2.S1 ----
The area-scaled violin plots show monotonic increases in the review rating score as trait scores rise from 0 to 1, with the largest effects for patient-oriented subjective judgments.
%% {CC-content-v0531}: restates the figure caption (monotonic rise; patient-oriented judgments). | cut or compress to a pointer to Fig.~\ref{fig:trait_satisfaction}. ========>
%
%% ---- P2.S2 ----
Sensitivity to Patient Satisfaction (SPS) had the strongest correlation ($r = 0.814$, 95\% CI [0.812, 0.815], $n = 186{,}250$), followed by Interpersonal Qualities \& Communication ($r = 0.804$ [0.802, 0.806]) and Social Trust Signals ($r = 0.789$ [0.787, 0.790]); all $p<0.001$.
%
%% ---- P2.S3 [v0519-rev3: drop "losers list" + sentiment->rating; emphasize non-zero residuals across all 10] ----
Partial correlations controlling for review rating scores remain non-zero across all ten traits, with Social Trust Signals ($r_{\text{partial}} = 0.28$) and Conscientiousness ($0.21$) retaining the most independent signal (attenuation range 66--96\%).
%
%% ---- P2.S3b ----
All trait--satisfaction CIs are 1{,}000-resample bootstrap 95\% intervals.


% =========================================================
% Para [trait-rating.regression] Result -- Big Five traits independently predict the rating (R^2 = 0.803)
% =========================================================
%
%% ---- P3.S1 ----
Among the Big Five, Agreeableness ($r = 0.747$, 95\% CI [0.745, 0.749]) and Conscientiousness ($r = 0.693$ [0.691, 0.696]) had the strongest associations with the review rating score, and Extraversion the weakest ($r = 0.412$ [0.407, 0.416]).
%
%% ---- P3.S2 ----
Violin shapes tighten at higher trait levels: physicians scoring 1.0 cluster around 4--5 stars while those scoring 0 spread broadly toward lower ratings.
%% {CC-content-v0531}: this distribution-shape aside interrupts the Big-Five -> regression flow (S1 then S3). | move next to the violin discussion in P2, or after S3. ========>
%
%% ---- P3.S3 ----
A multiple regression of standardized trait scores on the review rating score (10 predictors, $n = 83{,}230$, $R^2 = 0.803$) showed that these traits independently contribute to the review rating score.
%
%% ---- P3.S4 ----
Social Trust Signals was the strongest standardized predictor ($\beta = 0.213$, 95\% CI [0.207, 0.219]), followed by Interpersonal Qualities \& Communication ($\beta = 0.134$ [0.124, 0.144]) and Conscientiousness ($\beta = 0.132$ [0.128, 0.136]).
%
%% ---- P3.S4b ----
All 10 predictors' CIs are bounded away from zero, down to Emotional Stability ($\beta = 0.015$, 95\% CI [0.010, 0.020]).
%
%% ---- P3.S5 [v0519-rev3: synthesis -- mirrors abstract P1.S4c; affirmative framing] ----
From the above analysis, trait scores correlate strongly with review rating scores yet each retains substantial independent variance.
%
